\chapter{Memorias Digitales}

\section{¿Qué es una memoria?}

La memoria es la parte de un sistema digital que almacena datos binarios en grandes cantidades. Está formada por matrices de elementos de almacenamiento ---biestables o capacitores--- capaces de retener un bit de información cada uno.

En un sistema microprogramable la memoria cumple dos funciones distintas:

\begin{itemize}
  \item \textbf{Memoria de programa:} almacena las instrucciones que la CPU va a ejecutar. No se modifica durante la ejecución del programa. Se implementa en memoria no volátil (Flash, ROM).
  \item \textbf{Memoria de datos:} almacena las variables y resultados intermedios del programa. Se modifica constantemente durante la ejecución. Se implementa en RAM (volátil).
\end{itemize}

\section{Unidades de datos binarios}

\begin{table}[H]
\centering
\begin{tabular}{llp{7cm}}
\toprule
Unidad & Tamaño & Descripción \\
\midrule
Bit & 1 bit & Unidad mínima de información: un 0 o un 1 \\
Nibble & 4 bits & Medio byte; útil en representación BCD \\
Byte & 8 bits & Unidad básica de almacenamiento en la mayoría de los sistemas \\
Palabra (\textit{word}) & variable & Unidad completa de trabajo de la CPU (8, 16, 32 o 64 bits según la arquitectura) \\
\bottomrule
\end{tabular}
\end{table}

Las capacidades de memoria se expresan en múltiplos de byte:

\begin{table}[H]
\centering
\begin{tabular}{ll}
\toprule
Múltiplo & Equivale a \\
\midrule
1 Kbyte (KB) & $2^{10}$ bytes = 1.024 bytes \\
1 Megabyte (MB) & $2^{20}$ bytes = 1.048.576 bytes \\
1 Gigabyte (GB) & $2^{30}$ bytes = 1.073.741.824 bytes \\
\bottomrule
\end{tabular}
\end{table}

\section{Organización interna: la matriz de memoria}

Internamente, una memoria se organiza como una \textbf{matriz de celdas}, donde cada celda almacena un bit. La forma en que se agrupan esas celdas determina la organización de la memoria.

\imagenPendiente{matriz de celdas --- tres organizaciones posibles (64×1, 8×8, 16×4)}

Por ejemplo, una matriz de 64 celdas puede organizarse como:

\begin{itemize}
  \item \textbf{64 $\times$ 1:} 64 palabras de 1 bit cada una.
  \item \textbf{8 $\times$ 8:} 8 palabras de 8 bits (8 bytes).
  \item \textbf{16 $\times$ 4:} 16 palabras de 4 bits (16 nibbles).
\end{itemize}

La \textbf{capacidad} de una memoria se calcula como:

$$\text{Capacidad} = N_{\text{direcciones}} \times \text{longitud de palabra}$$

\textbf{Ejemplos:}

\begin{itemize}
  \item Una memoria 16 $\times$ 4: 16 posiciones $\times$ 4 bits = 64 bits en total.
  \item Una memoria 16K $\times$ 8: $16.384$ posiciones $\times$ 8 bits = 131.072 bits = 16 KB.
\end{itemize}

\begin{notabox}
Los números de posiciones son siempre potencias de 2. El ``16k'' es en realidad $2^{14} = 16.384$, redondeado a la unidad más próxima por convención.
\end{notabox}

\section{Direccionamiento}

Cada posición de la matriz tiene una \textbf{dirección}: un código binario único que la identifica. Para acceder a una posición se coloca su dirección en el \textbf{bus de direcciones}; un decodificador interno selecciona la fila y columna correspondientes.

\imagenPendiente{esquema de una memoria con bus de direcciones, bus de datos y señales de control (CS, OE, WE)}

La cantidad de posiciones direccionables depende del número de líneas del bus de direcciones:

$$N_{\text{posiciones}} = 2^{N_{\text{líneas de dirección}}}$$

\textbf{Ejemplo:} una memoria con 10 líneas de dirección (A0--A9) puede tener $2^{10} = 1.024$ posiciones.

\section{Operaciones básicas: lectura y escritura}

Para controlar si una operación es de lectura o escritura, y para habilitar o deshabilitar el chip, las memorias disponen de señales de control. Las más comunes son:

\begin{table}[H]
\centering
\begin{tabular}{llp{7cm}}
\toprule
Señal & Nombre & Función \\
\midrule
$\overline{CS}$ (o $\overline{CE}$) & \textit{Chip Select / Chip Enable} & Activa el chip (activa en bajo). Permite conectar varios chips al mismo bus y seleccionar solo uno a la vez \\
$\overline{OE}$ & \textit{Output Enable} & Habilita las salidas de datos durante una lectura \\
$\overline{WE}$ & \textit{Write Enable} & Habilita la escritura de datos \\
\bottomrule
\end{tabular}
\end{table}

\subsection{Operación de escritura}

\begin{enumerate}
  \item Se coloca la dirección en el bus de direcciones.
  \item Se activa $\overline{CS}$ (se pone en bajo).
  \item Se coloca el dato en el bus de datos.
  \item Se activa $\overline{WE}$ (pulso en bajo): el dato queda grabado en la posición seleccionada.
\end{enumerate}

\imagenPendiente{diagrama de tiempos de una operación de escritura}

\begin{advertenciabox}
Escribir en una posición \textbf{destruye} el dato que había antes en esa dirección.
\end{advertenciabox}

\subsection{Operación de lectura}

\begin{enumerate}
  \item Se coloca la dirección en el bus de direcciones.
  \item Se activa $\overline{CS}$ (se pone en bajo).
  \item Se activa $\overline{OE}$ (se pone en bajo): la memoria deposita el dato almacenado en esa dirección sobre el bus de datos.
  \item La CPU lee el dato del bus.
\end{enumerate}

\imagenPendiente{diagrama de tiempos de una operación de lectura}

\begin{notabox}
La lectura es \textbf{no destructiva}: el dato permanece almacenado en la memoria después de ser leído.
\end{notabox}

\section{Clasificación de las memorias}

Las memorias se clasifican según dos criterios principales:

\textbf{Según su posibilidad de escritura:}
\begin{itemize}
  \item \textbf{ROM} (\textit{Read Only Memory}): solo lectura.
  \item \textbf{RAM} (\textit{Random Access Memory}): lectura y escritura.
\end{itemize}

\textbf{Según su volatilidad:}
\begin{itemize}
  \item \textbf{Volátiles:} pierden los datos al retirar la alimentación (RAM).
  \item \textbf{No volátiles:} conservan los datos sin alimentación (ROM y sus variantes).
\end{itemize}

\imagenPendiente{árbol de clasificación de memorias (RAM $\rightarrow$ SRAM / DRAM; ROM $\rightarrow$ máscara / PROM / EPROM / EEPROM / Flash)}

\section{Memorias RAM}

\subsection{RAM estática (SRAM)}

Utiliza \textbf{biestables (latches o flip-flops)} como elementos de almacenamiento. Cada celda almacena un bit de forma estable mientras haya alimentación, sin necesidad de refresco.

\imagenPendiente{celda básica SRAM (latch con líneas de selección de bit y datos)}

\textbf{Características:}
\begin{itemize}
  \item Mantiene el dato indefinidamente con alimentación.
  \item Alta velocidad de acceso.
  \item Mayor área de silicio por bit $\rightarrow$ mayor costo y menor densidad.
  \item Se usa en cachés de CPU, registros de microcontroladores y en la memoria de datos (SRAM) de los microcontroladores.
\end{itemize}

\subsection{RAM dinámica (DRAM)}

Utiliza \textbf{capacitores} como elementos de almacenamiento. Un transistor controla el acceso al capacitor. La carga del capacitor representa el bit almacenado.

\imagenPendiente{celda básica DRAM (transistor MOS + capacitor)}

\textbf{Características:}
\begin{itemize}
  \item El capacitor se descarga lentamente $\rightarrow$ necesita un \textbf{ciclo de refresco} periódico para no perder el dato.
  \item Mayor densidad de almacenamiento (celda más pequeña) $\rightarrow$ menor costo por bit.
  \item Requiere circuitería adicional para el refresco.
  \item Es la tecnología utilizada en los módulos de RAM de las computadoras personales (DDR, DDR4, etc.).
\end{itemize}

\section{Memorias ROM y sus variantes}

\subsection{ROM de máscara}

Programada durante la fabricación del chip. El contenido no puede modificarse de ninguna manera. Solo es económica en grandes volúmenes de producción.

\subsection{PROM (\textit{Programmable ROM})}

Puede ser programada \textbf{una única vez} por el usuario mediante un programador de memoria. El proceso es irreversible: una vez programada no puede cambiarse. También llamada OTP (\textit{One Time Programmable}).

\subsection{EPROM (\textit{Erasable PROM})}

Puede ser \textbf{borrada mediante luz ultravioleta} a través de una ventana de cuarzo en el encapsulado, y luego reprogramada. El borrado es total (no selectivo) y tarda varios minutos. Reconocible por la ventana transparente en el chip.

\imagenPendiente{foto o dibujo de un chip EPROM con ventana de cuarzo}

\subsection{EEPROM (\textit{Electrically Erasable PROM})}

Puede ser \textbf{borrada y reprogramada eléctricamente}, byte a byte, sin necesidad de extraer el chip del circuito. Permite actualizaciones en el propio sistema (\textit{In-System Programming}). Velocidad de escritura más lenta que la RAM; número de ciclos de escritura limitado (típicamente 100.000 a 1.000.000).

\subsection{Memoria Flash}

Es la evolución de la EEPROM. Comparte la capacidad de borrado y escritura eléctrica, pero opera por \textbf{bloques} en lugar de byte a byte, lo que la hace significativamente más rápida en operaciones de borrado. Combina:

\begin{itemize}
  \item Alta densidad de almacenamiento.
  \item No volatilidad.
  \item Lectura/escritura en el propio circuito.
  \item Velocidad de operación comparativamente alta.
  \item Bajo costo por bit.
\end{itemize}

Por estas razones es la tecnología dominante en memorias USB, tarjetas SD y en la \textbf{memoria de programa de los microcontroladores modernos} (incluyendo el ATmega328, que tiene 32 KB de Flash).

\textbf{Comparativa de tipos de ROM:}

\begin{table}[H]
\centering
\begin{tabular}{llll}
\toprule
Tipo & Programación & Borrado & Granularidad borrado \\
\midrule
ROM de máscara & En fábrica & No borrable & --- \\
PROM & Usuario, 1 vez & No borrable & --- \\
EPROM & Usuario & UV, total & Todo el chip \\
EEPROM & Eléctricamente & Eléctricamente & Byte a byte \\
Flash & Eléctricamente & Eléctricamente & Por bloques \\
\bottomrule
\end{tabular}
\end{table}

\section{Expansión de memoria}

En muchas aplicaciones la capacidad o el ancho de palabra de un único chip de memoria no son suficientes. Se pueden combinar varios chips para ampliar la memoria total del sistema. Existen dos técnicas:

\subsection{Expansión de palabra (ampliación de bits por posición)}

Se usa cuando el ancho de palabra de los chips es menor que el del bus de datos del sistema. Los chips se conectan \textbf{en paralelo}, compartiendo el mismo bus de direcciones y las mismas señales de control, pero cada uno aporta algunos bits del bus de datos.

\textbf{Ejemplo:} el bus de datos del sistema es de 8 bits, pero se tienen chips de 4 bits de ancho (organizados como 1K $\times$ 4). Se conectan dos chips en paralelo: uno proporciona D0--D3 y el otro D4--D7. Ambos chips responden a la misma dirección simultáneamente.

\imagenPendiente{diagrama de expansión de palabra --- dos chips 1K×4 formando una memoria 1K×8}

$$\text{Chips necesarios} = \frac{\text{ancho de palabra del sistema}}{\text{ancho de palabra del chip}}$$

\subsection{Expansión de capacidad (ampliación de direcciones)}

Se usa cuando la cantidad de posiciones de un único chip es menor que la requerida. Los chips se conectan de forma que cada uno responde a un rango diferente de direcciones. Se usa la señal $\overline{CS}$ de cada chip para habilitarlo solo cuando el bus de direcciones está en su rango.

\textbf{Ejemplo:} se necesitan 4K posiciones de 8 bits, pero solo se tienen chips de 1K $\times$ 8. Se conectan 4 chips: el primero responde a las direcciones 0x000--0x3FF, el segundo a 0x400--0x7FF, y así sucesivamente. Un decodificador (por ejemplo, un 74138) genera las señales $\overline{CS}$ individuales a partir de las líneas de mayor peso del bus de direcciones.

\imagenPendiente{diagrama de expansión de capacidad --- cuatro chips 1K×8 formando una memoria 4K×8, con decodificador de dirección}

$$\text{Chips necesarios} = \frac{\text{capacidad total requerida}}{\text{capacidad de cada chip}}$$

\subsection{Expansión combinada}

Es posible aplicar ambas técnicas simultáneamente para ampliar tanto el ancho de palabra como la cantidad de posiciones. En ese caso:

$$\text{Chips necesarios} = \frac{\text{capacidad total}}{\text{capacidad del chip}} \times \frac{\text{ancho de palabra del sistema}}{\text{ancho de palabra del chip}}$$
